\chapter{Campos y espacios vectoriales}
\label{cap:campos-espacios}

El álgebra lineal estudia conjuntos con operaciones que permiten sumar vectores y multiplicarlos por escalares. Nos centraremos en espacios vectoriales de dimensión finita: aquellos que tienen una base con un número finito de vectores. El conjunto de vectores puede tener infinitos elementos; por ejemplo, $\mathbb{R}^2$ sí tiene dimensión finita. Estas estructuras dan un marco común para estudiar matrices, sistemas de ecuaciones lineales y transformaciones lineales.
\section{Campos y operaciones algebraicas}
Antes de introducir los espacios vectoriales debemos precisar de dónde provienen sus escalares. Trabajaremos principalmente con los racionales $\mathbb{Q}$, los reales $\mathbb{R}$ y los complejos $\mathbb{C}$. Los tres son ejemplos de campos: conjuntos en los que se puede sumar, restar, multiplicar y dividir por cualquier elemento distinto de cero, conservando las reglas algebraicas usuales.
\begin{definition}
Sea $K$ un conjunto no vacío provisto de dos operaciones
\(+\colon K\times K\to K\) y \(\cdot\colon K\times K\to K\).
Diremos que $(K,+,\cdot)$ es un \indexemph{campo} si para cualesquiera $a,b,c\in K$ se satisfacen las siguientes propiedades:
\begin{enumerate}
\item La suma es asociativa y conmutativa: $(a+b)+c=a+(b+c)$ y $a+b=b+a$.
\item Existe un neutro aditivo $0\in K$, con $a+0=a$, y cada $a\in K$ tiene un inverso $-a\in K$, con $a+(-a)=0$.
\item El producto es asociativo y conmutativo: $(ab)c=a(bc)$ y $ab=ba$.
\item Existe un neutro multiplicativo $1\in K$, con $1\neq0$ y $a1=a$. Cada $a\neq0$ tiene un inverso multiplicativo $a^{-1}$, con $aa^{-1}=1$.
\item El producto distribuye sobre la suma: $a(b+c)=ab+ac$.
\end{enumerate}
\end{definition}
\noindent En particular, la suma hace de $K$ un grupo abeliano, mientras que los elementos no nulos forman un grupo abeliano bajo el producto. Las aplicaciones $+$ y $\cdot$ ya incorporan el cierre de las operaciones en $K$.
\begin{example}
El conjunto $\mathbb{Q}$ de los racionales, con las operaciones usuales, es un campo: si $a\in\mathbb{Q}$ y $a\neq0$, entonces $a^{-1}=1/a\in\mathbb{Q}$.
\end{example}
\begin{example}
Los conjuntos $\mathbb{R}$ y $\mathbb{C}$, con sus operaciones usuales, también son campos y serán nuestros escalares principales. Cuando una afirmación valga para ambos, escribiremos $K=\mathbb{R}$ o $K=\mathbb{C}$.
\end{example}
\begin{example}
Si $p$ es primo, el conjunto $\mathbb{F}_p=\mathbb{Z}/p\mathbb{Z}$, con suma y producto módulo $p$, es un campo. Por ejemplo, $\mathbb{F}_2=\{0,1\}$ y $1+1=0$, mientras que $1\cdot1=1$. Así, un campo también puede tener un número finito de elementos.
\end{example}
\begin{example}
Los enteros $\mathbb{Z}$ no forman un campo: aunque tienen las operaciones y los neutros usuales, $2$ no tiene inverso multiplicativo en $\mathbb{Z}$; no existe $n\in\mathbb{Z}$ tal que $2n=1$.
\end{example}
A partir de los axiomas de campo pueden deducirse numerosas propiedades que utilizamos de manera habitual en manipulaciones algebraicas.
\begin{proposition}
Sea $K$ un campo. Entonces se cumplen las siguientes propiedades:
\begin{enumerate}
\item[a)] El elemento neutro aditivo es único.
\item[b)] Para cada $a\in K$, el inverso aditivo $-a$ es único.
\item[c)] El elemento neutro multiplicativo es único.
\item[d)] Para cada $a\in K\setminus\{0\}$, el inverso multiplicativo $a^{-1}$ es único.
\end{enumerate}
\end{proposition}
\begin{proof}
Si $0$ y $0'$ son neutros aditivos, entonces $0=0+0'=0'$. Si $b$ y $c$ son inversos aditivos de $a$, se tiene $b=b+0=b+(a+c)=(b+a)+c=0+c=c$. Los mismos argumentos multiplicativos dan $1=1\cdot1'=1'$ y, si $ab=ac=1$ con $a\neq0$, entonces $b=b\cdot1=b(ac)=(ba)c=1\cdot c=c$. En cada caso, el elemento con la propiedad requerida es único.
\end{proof}
\begin{example}
En $\mathbb{R}$, los inversos de $5$ son $-5$ y $5^{-1}=1/5$, pues $5+(-5)=0$ y $5(1/5)=1$. La proposición garantiza que no hay otros elementos con esas propiedades.
\end{example}
Una herramienta fundamental para resolver ecuaciones y realizar manipulaciones algebraicas es la ley de cancelación.
\begin{proposition}[Ley de cancelación]
Sea $K$ un campo y sean $a,b,c\in K$. Entonces:
\begin{enumerate}
\item[a)] Si $a+c=b+c$, entonces $a=b$.
\item[b)] Si $c\neq0$ y $ac=bc$, entonces $a=b$.
\end{enumerate}
\end{proposition}
\begin{proof}
En la primera igualdad sumamos $-c$ a ambos lados y obtenemos $a+c-c=b+c-c$, de donde $a=b$. En la segunda multiplicamos por $c^{-1}$: $acc^{-1}=bcc^{-1}$, así que también $a=b$.
\end{proof}
Las siguientes identidades, aunque elementales, aparecerán continuamente al trabajar con espacios vectoriales, matrices y transformaciones lineales.
\begin{proposition}
Sea $K$ un campo y sean $a,b\in K$. Entonces:
\begin{enumerate}
\item[a)] $a\cdot0=0$.
\item[b)] $(-a)b=a(-b)=-(ab)$.
\item[c)] $(-a)(-b)=ab$.
\item[d)] $-(-a)=a$.
\item[e)] $(-1)a=-a$.
\end{enumerate}
\end{proposition}
\begin{proof}
Por distributividad, $a\cdot0=a(0+0)=a\cdot0+a\cdot0$; al cancelar un sumando resulta $a\cdot0=0$. También $ab+(-a)b=(a-a)b=0$, de modo que $(-a)b=-(ab)$; análogamente, $a(-b)=-(ab)$. Aplicar estas identidades dos veces da $(-a)(-b)=ab$. La unicidad del inverso aditivo implica $-(-a)=a$ y, como $a+(-1)a=(1-1)a=0$, también $(-1)a=-a$.
\end{proof}
\begin{proposition}[Propiedad del producto nulo]
Sean $a,b\in K$. Si $ab=0$, entonces $a=0$ o $b=0$.
\end{proposition}
\begin{proof}
Si $a=0$, no hay nada que demostrar. Si $a\neq0$, multiplicamos $ab=0$ por $a^{-1}$ y obtenemos $b=a^{-1}0=0$.
\end{proof}

\section{Los campos \texorpdfstring{$\mathbb{R}$}{R} y \texorpdfstring{$\mathbb{C}$}{C}}
Al hablar de un espacio vectorial siempre especificaremos su campo de escalares. En lo que sigue, $K$ denotará principalmente a $\mathbb{R}$ o a $\mathbb{C}$. Los complejos amplían a los reales al incorporar una solución de $x^2=-1$; al mismo tiempo, se pueden representar como puntos del plano. Esta doble lectura, algebraica y geométrica, será útil para entender las operaciones antes de generalizarlas a $K^n$.
\begin{definition}
El conjunto de los \indexemph{números complejos} es
$\mathbb{C}=\{a+bi:a,b\in\mathbb{R}\}$, donde $i^2=-1$.
Si $z=a+bi$, su \indexemph{parte real} y su \indexemph{parte imaginaria} son
$\operatorname{Re}(z)=a$ e $\operatorname{Im}(z)=b$. La igualdad $a+bi=c+di$
equivale a que $a=c$ y $b=d$.
\end{definition}
Identificamos $a+bi$ con el punto $(a,b)\in\mathbb{R}^2$. Esta correspondencia
preserva la suma y la multiplicación por escalares reales, así que identifica
$\mathbb{C}$ con $\mathbb{R}^2$ como espacios vectoriales sobre $\mathbb{R}$.
El producto complejo, en cambio, tiene la regla propia
$(a,b)(c,d)=(ac-bd,ad+bc)$.
\begin{example}
Si $z=3+2i$, entonces corresponde al punto $(3,2)$ y sus coordenadas son
$\operatorname{Re}(z)=3$ e $\operatorname{Im}(z)=2$. De igual modo, $-4-7i$
corresponde a $(-4,-7)$.
\end{example}
\begin{example}
Todo real es un complejo cuya parte imaginaria es cero: $a=a+0i$, de modo que
$\mathbb{R}\subseteq\mathbb{C}$. Por ejemplo, $5$ se representa como $(5,0)$.
\end{example}
La identificación con pares también hace visibles las operaciones. Al multiplicar,
la regla resulta de distribuir y sustituir $i^2$ por $-1$.
\begin{definition}[Aritmética de los números complejos]
Si $z=a+bi$ y $w=c+di$, entonces
$z+w=(a+c)+(b+d)i$ y $zw=(ac-bd)+(ad+bc)i$.
En coordenadas, esto equivale a
$(a,b)+(c,d)=(a+c,b+d)$ y $(a,b)(c,d)=(ac-bd,ad+bc)$.
\end{definition}
\begin{example}
Para $z=2+3i$ y $w=4-i$, se obtiene $z+w=6+2i$ y
$zw=(2+3i)(4-i)=8-2i+12i-3i^2=11+10i$.
\end{example}
\begin{figure}[htbp]
  \centering
  \begin{tikzpicture}[
  >=Latex,
  font=\small,
  line cap=round,
  line join=round
]
  \draw[->,gray] (-0.35,0) -- (4.0,0) node[right] {$\operatorname{Re}$};
  \draw[->,gray] (0,-0.3) -- (0,3.65) node[above] {$\operatorname{Im}$};
  \draw[densely dashed,gray] (2,1) -- (3,3);
  \draw[densely dashed,gray] (1,2) -- (3,3);
  \draw[->,blue!70!black,very thick] (0,0) -- (2,1)
    node[below right] {$z=2+i$};
  \draw[->,orange!85!black,very thick] (0,0) -- (1,2)
    node[above left] {$w=1+2i$};
  \draw[->,blue!70!black,very thick] (1,2) -- (3,3);
  \draw[->,orange!85!black,very thick] (2,1) -- (3,3);
  \fill[black] (3,3) circle (1.6pt)
    node[above right] {$z+w=3+3i$};
\end{tikzpicture}

  \caption{La suma compleja se obtiene sumando las componentes, o trasladando uno de los vectores para completar el paralelogramo.}
  \label{fig:suma-compleja}
\end{figure}
\begin{example}
Como $i^2=-1$, las potencias se repiten cada cuatro términos:
$i^{4k}=1$, $i^{4k+1}=i$, $i^{4k+2}=-1$ e $i^{4k+3}=-i$.
\end{example}
\begin{example}
Para $z=1+2i$, tenemos $z^2=(1+2i)^2=1+4i+4i^2=-3+4i$.
\end{example}
Otra operación fundamental en $\mathbb{C}$ es la conjugación, la cual posee una interpretación geométrica sencilla: refleja un número complejo respecto del eje real en el plano complejo.
\begin{definition}
Si $z=a+bi\in\mathbb{C}$, su \indexemph{conjugado} es $\overline{z}=a-bi$.
\end{definition}
\begin{example}
Por ejemplo, $\overline{3+5i}=3-5i$ y $\overline{-2-7i}=-2+7i$.
En particular, $z\in\mathbb{R}$ si y sólo si $\overline{z}=z$.
\end{example}
\begin{proposition}
Sean $z,w\in\mathbb{C}$. Entonces:
\begin{enumerate}
\item[a)] $\overline{\overline{z}}=z$.
\item[b)] $\overline{z+w}=\overline{z}+\overline{w}$.
\item[c)] $\overline{zw}=\overline{z}\,\overline{w}$.
\item[d)] $\overline{z^{-1}}=(\overline{z})^{-1}$, siempre que $z\neq0$.
\item[e)] $z+\overline{z}=2\operatorname{Re}(z)$.
\item[f)] $z-\overline{z}=2i\operatorname{Im}(z)$.
\item[g)] $z\overline{z}=(\operatorname{Re}z)^2+(\operatorname{Im}z)^2$.
\end{enumerate}
\end{proposition}
\begin{proof}
Al escribir $z=a+bi$ y $w=c+di$, cada identidad se verifica sustituyendo las definiciones; por ejemplo,
$z\overline{z}=(a+bi)(a-bi)=a^2+b^2$.
\end{proof}
El conjugado permite definir la magnitud y la distancia en el plano complejo, usando la distancia euclidiana.
\begin{definition}
Para $z=a+bi$, su \indexemph{módulo} es $|z|=\sqrt{a^2+b^2}=\sqrt{z\overline{z}}$.
La distancia entre $z,w\in\mathbb{C}$ se define como $d(z,w)=|z-w|$; por la
identificación con $\mathbb{R}^2$, es la distancia euclidiana entre los puntos.
\end{definition}
\begin{proposition}
Para cualesquiera $z,w\in\mathbb{C}$ se cumplen
$|z|^2=z\overline{z}$, $|zw|=|z||w|$, $|\overline{z}|=|z|$ y
$|z+w|\leq |z|+|w|$. Además, si $z\neq0$, entonces
$z^{-1}=\overline{z}/|z|^2$.
\end{proposition}
\begin{proof}
De $z\overline{z}=|z|^2$ y $\overline{zw}=\overline z\,\overline w$ se obtiene
$|zw|^2=|z|^2|w|^2$, así que $|zw|=|z||w|$; la igualdad
$|\overline z|=|z|$ se sigue del mismo cálculo. Para la desigualdad triangular,
escribimos $|z+w|^2=|z|^2+|w|^2+2\operatorname{Re}(z\overline w)$ y usamos
$\operatorname{Re}(z\overline w)\leq|z||w|$. Finalmente,
$z\,\overline z/|z|^2=1$, por lo que $z^{-1}=\overline z/|z|^2$.
\end{proof}
\begin{example}
Para $z=3+4i$, el módulo es $|z|=5$. Si $z=1+2i$ y $w=4+6i$, entonces
$d(z,w)=|z-w|=|-3-4i|=5$. Finalmente, para $z=2+3i$,
$z^{-1}=(2-3i)/13$.
\end{example}
\begin{figure}[htbp]
  \centering
  \begin{tikzpicture}[
  >=Latex,
  font=\small,
  line cap=round,
  line join=round
]
  \draw[->,gray] (-0.45,0) -- (3.5,0) node[right] {$\operatorname{Re}$};
  \draw[->,gray] (0,-2.0) -- (0,2.5) node[above] {$\operatorname{Im}$};
  \draw[densely dashed,gray] (2,1) -- (2,0);
  \draw[densely dashed,gray] (2,-1) -- (2,0);
  \draw[->,blue!70!black,very thick] (0,0) -- (2,1)
    node[above right] {$z=2+i$};
  \draw[->,orange!85!black,very thick] (0,0) -- (2,-1)
    node[below right] {$\overline z=2-i$};
  \draw[<->,red!70!black,thick] (2.3,-1) -- (2.3,1)
    node[midway,right] {$2|\operatorname{Im}z|$};
  \node[below left] at (0,0) {$0$};
  \node[above left] at (1.0,0.5) {$|z|=\sqrt{5}$};
\end{tikzpicture}

  \caption{El conjugado refleja $z$ respecto del eje real; el módulo es la distancia al origen.}
  \label{fig:conjugado-modulo}
\end{figure}
\begin{figure}[htbp]
  \centering
  \begin{tikzpicture}[
  >=Latex,
  font=\small,
  line cap=round,
  line join=round
]
  \draw[->,gray] (-0.35,0) -- (3.2,0) node[right] {$\operatorname{Re}$};
  \draw[->,gray] (0,-0.35) -- (0,2.8) node[above] {$\operatorname{Im}$};
  \draw[densely dashed,gray] (2,1) -- (2,0);
  \draw[densely dashed,gray] (1,2) -- (1,0);
  \fill[blue!70!black] (2,1) circle (1.7pt)
    node[below right] {$z=2+i$};
  \fill[orange!85!black] (1,2) circle (1.7pt)
    node[above left] {$w=1+2i$};
  \draw[<->,red!70!black,very thick] (2,1) -- (1,2)
    node[midway,above right] {$d(z,w)=|z-w|=\sqrt{2}$};
\end{tikzpicture}

  \caption{La distancia entre dos complejos es la longitud del segmento que une sus puntos en el plano.}
  \label{fig:distancia-complejos}
\end{figure}
La representación cartesiana $z=a+bi$ también puede expresarse mediante la distancia al origen y el ángulo respecto del eje real positivo.
\begin{definition}
Sea $z=a+bi\neq0$. Un número $\theta\in\mathbb{R}$ es un \indexemph{argumento}
de $z$ si $a=|z|\cos\theta$ y $b=|z|\sin\theta$. Si $\theta$ es uno de ellos,
el conjunto de todos se denota por
$\arg(z)=\{\theta+2k\pi:k\in\mathbb{Z}\}$. El único argumento en $(-\pi,\pi]$
se llama \indexemph{argumento principal} y se denota por $\operatorname{Arg}(z)$.
\end{definition}
La notación $\arg(z)$ representa una familia de ángulos, mientras que
$\operatorname{Arg}(z)$ elige uno. Esa elección es continua salvo al cruzar el
eje real negativo, donde pasa de valores cercanos a $\pi$ a valores cercanos a
$-\pi$; a este salto se le asocia el corte de la rama principal.
\begin{figure}[htbp]
  \centering
  \begin{tikzpicture}[
  >=Latex,
  font=\small,
  line cap=round,
  line join=round
]
  \draw[->,gray] (-3.5,0) -- (2.0,0) node[right] {$\operatorname{Re}$};
  \draw[->,gray] (0,-1.7) -- (0,1.7) node[above] {$\operatorname{Im}$};
  \draw[red!70!black,dashed,very thick] (-3.25,0) -- (0,0)
    node[midway,below] {corte de rama};
  \fill[blue!70!black] (-2.4,0.55) circle (1.7pt)
    node[above left] {$z_+$};
  \fill[orange!85!black] (-2.4,-0.55) circle (1.7pt)
    node[below left] {$z_-$};
  \draw[->,blue!70!black,thick] (0.65,0) arc[start angle=0,end angle=167,radius=0.65]
    node[midway,above right] {$\operatorname{Arg}(z_+)\to\pi$};
  \draw[->,orange!85!black,thick] (0.95,0) arc[start angle=0,end angle=-167,radius=0.95]
    node[midway,below right] {$\operatorname{Arg}(z_-)\to-\pi$};
  \node[align=center] at (1.0,1.45) {$\operatorname{Arg}(z)\in(-\pi,\pi]$\\
    al cruzar el eje real negativo\\la rama salta en $2\pi$};
\end{tikzpicture}

  \caption{La rama principal del argumento toma valores en $(-\pi,\pi]$ y presenta un salto a lo largo del eje real negativo.}
  \label{fig:argumento-principal}
\end{figure}
\begin{example}
Para $z=1+i$, se tiene $|z|=\sqrt2$ y un argumento es $\theta=\pi/4$.
Por ello, $\arg(z)=\{\pi/4+2k\pi:k\in\mathbb{Z}\}$ y
$\operatorname{Arg}(z)=\pi/4$.
\end{example}
Si $\theta$ es un argumento de $z\neq0$, entonces
$z=|z|(\cos\theta+i\sin\theta)=|z|e^{i\theta}$, usando la identidad de Euler
$e^{i\theta}=\cos\theta+i\sin\theta$. Esta es la \indexemph{forma polar} de $z$.
\begin{proposition}
Sean $z=r(\cos\theta+i\sin\theta)$ y $w=s(\cos\varphi+i\sin\varphi)$, con
$r,s\geq0$. Entonces
$zw=rs\bigl(\cos(\theta+\varphi)+i\sin(\theta+\varphi)\bigr)$.
En particular, $|zw|=|z||w|$ y, si $z,w\neq0$,
$\arg(zw)=\arg(z)+\arg(w)\pmod{2\pi}$.
\end{proposition}
\begin{proof}
Al distribuir el producto y aplicar las identidades trigonométricas de suma,
\[
\begin{aligned}
zw
&=rs(\cos\theta+i\sin\theta)(\cos\varphi+i\sin\varphi)\\
&=rs\bigl((\cos\theta\cos\varphi-\sin\theta\sin\varphi)
+i(\sin\theta\cos\varphi+\cos\theta\sin\varphi)\bigr)\\
&=rs\bigl(\cos(\theta+\varphi)+i\sin(\theta+\varphi)\bigr).
\end{aligned}
\end{proof}
La forma polar permite calcular potencias de números complejos de manera particularmente sencilla.
\begin{proposition}
Si $z=r(\cos\theta+i\sin\theta)$ y $n\in\mathbb{N}$, entonces
$z^n=r^n\bigl(\cos(n\theta)+i\sin(n\theta)\bigr)$, o, equivalentemente,
$(re^{i\theta})^n=r^ne^{in\theta}$.
\end{proposition}
\begin{definition}
Sean $z\in\mathbb{C}$ y $n\geq1$. Un número $w\in\mathbb{C}$ es una
\indexemph{raíz $n$-ésima} de $z$ si $w^n=z$. Para $z=0$, la única raíz es
$w=0$. Si $z=r(\cos\theta+i\sin\theta)$ con $r>0$, sus $n$ raíces son
\[
w_k=r^{1/n}
\left(
\cos\frac{\theta+2k\pi}{n}
+i\sin\frac{\theta+2k\pi}{n}
\right),
\qquad k=0,\ldots,n-1.
\]
\end{definition}
\begin{example}
Las raíces cuadradas de $1$ son $1$ y $-1$.
\end{example}
\begin{example}
Si $\omega=e^{2\pi i/3}$, las raíces cúbicas de $1$ son $1$, $\omega$ y
$\omega^2$, es decir, $1$, $-\frac12+\frac{\sqrt3}{2}i$ y
$-\frac12-\frac{\sqrt3}{2}i$. En el plano son vértices de un triángulo
equilátero sobre la circunferencia unitaria.
Geométricamente, estos puntos son los vértices de un triángulo equilátero sobre la circunferencia unitaria.
\end{example}
\begin{figure}[htbp]
  \centering
  \begin{tikzpicture}[
  >=Latex,
  font=\small,
  line cap=round,
  line join=round
]
  \draw[->,gray] (-1.55,0) -- (1.65,0) node[right] {$\operatorname{Re}$};
  \draw[->,gray] (0,-1.45) -- (0,1.45) node[above] {$\operatorname{Im}$};
  \draw[gray!60] (0,0) circle (1);
  \foreach \ang/\lab/\pos in {0/$1$/right,120/$\omega$/above left,240/$\omega^2$/below left} {
    \draw[->,blue!65!black,thick] (0,0) -- (\ang:1);
    \fill[blue!65!black] (\ang:1) circle (1.6pt);
    \node[\pos] at (\ang:1.14) {\lab};
  }
  \node[align=center] at (0,-1.85) {$\omega^3=1$\\tres raíces sobre la circunferencia unitaria};
\end{tikzpicture}

  \caption{Las tres raíces cúbicas de la unidad se distribuyen con ángulos iguales sobre la circunferencia unitaria.}
  \label{fig:raices-unidad}
\end{figure}
\begin{theorem}[Fórmula de De Moivre]
Para todo $\theta\in\mathbb{R}$ y $n\in\mathbb{Z}$,
$(\cos\theta+i\sin\theta)^n=\cos(n\theta)+i\sin(n\theta)$.
Más generalmente, si $z=r(\cos\theta+i\sin\theta)$ con $r>0$, entonces
$z^n=r^n\bigl(\cos(n\theta)+i\sin(n\theta)\bigr)$.
\end{theorem}
\begin{proof}
Para $n\geq0$, aplicamos repetidamente la fórmula del producto en forma polar.
Si $n=-m<0$, usamos
$(\cos\theta+i\sin\theta)^{-1}=\cos(-\theta)+i\sin(-\theta)$ y aplicamos el
caso positivo a la potencia $m$.
\end{proof}
\begin{example}
Como $1+i=\sqrt2\,e^{i\pi/4}$, De Moivre da $(1+i)^8=(\sqrt2)^8e^{2\pi i}=16$.
\end{example}
\begin{example}
Como $|\sqrt3+i|=2$ y $\operatorname{Arg}(\sqrt3+i)=\pi/6$, resulta
$(\sqrt3+i)^3=8e^{i\pi/2}=8i$.
\end{example}
Los campos $\mathbb{R}$ y $\mathbb{C}$ sirven como escalares. A partir de cualquier
campo $K$ obtenemos los espacios coordenados $K^n$, primeros ejemplos de espacios
vectoriales de dimensión finita.
\begin{definition}
Sea $K$ un campo y $n\in\mathbb{N}$ con $n\geq1$. El conjunto de tuplas
$K^n=\{(x_1,\ldots,x_n):x_j\in K\}$ se opera componente a componente:
$(x_1,\ldots,x_n)+(y_1,\ldots,y_n)=(x_1+y_1,\ldots,x_n+y_n)$ y, para
$\lambda\in K$, $\lambda(x_1,\ldots,x_n)=(\lambda x_1,\ldots,\lambda x_n)$.
\end{definition}
Con estas operaciones, $K^n$ es un espacio vectorial sobre $K$. Para $n>1$ sus
elementos no forman un campo con estas operaciones: la multiplicación por
escalares toma un escalar y un vector, $K\times K^n\to K^n$.
\begin{example}
En $\mathbb{R}^2$, con $x=(2,-1)$ e $y=(3,4)$, la suma es $x+y=(5,3)$ y el
escalar $2$ produce $2x=(4,-2)$. Geométricamente, son vectores del plano.
\end{example}
\begin{example}
En $\mathbb{R}^3$, para $x=(1,2,-1)$ e $y=(0,-3,4)$ se tiene
$x+y=(1,-1,3)$ y $-2x=(-2,-4,2)$. Ahora las tuplas se representan como
vectores en el espacio tridimensional.
\end{example}
\begin{figure}[htbp]
  \centering
  \begin{tikzpicture}[
  >=Latex,
  font=\small,
  line cap=round,
  line join=round
]
  % El mismo vector se representa en el plano y en el espacio.
  \begin{scope}
    \draw[->,gray] (-0.25,0) -- (3.2,0) node[right] {$x_1$};
    \draw[->,gray] (0,-0.25) -- (0,2.25) node[above] {$x_2$};
    \draw[densely dashed,gray] (2,1) -- (2,0);
    \draw[densely dashed,gray] (2,1) -- (0,1);
    \draw[->,blue!70!black,very thick] (0,0) -- (2,1)
      node[above right] {$x=(2,1)$};
    \fill[blue!70!black] (2,1) circle (1.5pt);
    \node[font=\bfseries] at (1.35,2.65) {$\mathbb{R}^2$};
    \node[align=center] at (1.35,-0.65) {punto o vector\\en el plano};
  \end{scope}

  \begin{scope}[
    xshift=6.4cm,
    x={(0.95cm,0cm)},
    y={(0.32cm,0.2cm)},
    z={(0cm,0.9cm)}
  ]
    \draw[->,gray] (-0.2,0,0) -- (3.0,0,0) node[right] {$x_1$};
    \draw[->,gray] (0,-0.2,0) -- (0,2.1,0) node[right] {$x_2$};
    \draw[->,gray] (0,0,-0.2) -- (0,0,2.7) node[above] {$x_3$};
    \draw[densely dashed,gray] (2,1,2) -- (2,1,0);
    \draw[densely dashed,gray] (2,1,0) -- (2,0,0);
    \draw[densely dashed,gray] (2,1,0) -- (0,1,0);
    \draw[->,blue!70!black,very thick] (0,0,0) -- (2,1,2)
      node[above right] {$y=(2,1,2)$};
    \fill[blue!70!black] (2,1,2) circle (1.5pt);
    \node[font=\bfseries] at (1.15,0,2.7) {$\mathbb{R}^3$};
    \node[align=center] at (1.15,-0.55,-0.45) {punto o vector\\en el espacio};
  \end{scope}
\end{tikzpicture}

  \caption{Las tuplas reales de dos y tres coordenadas se representan como puntos o vectores en el plano y en el espacio.}
  \label{fig:coordenadas-r2-r3}
\end{figure}
\begin{example}
En $\mathbb{C}^2$, sean $z=(1+i,2-i)$ y $w=(3,-i)$. Entonces
$z+w=(4+i,2-2i)$. Para $\lambda=1+i$,
$\lambda z=((1+i)^2,(1+i)(2-i))=(2i,3+i)$.
\end{example}
\begin{example}
En general, si
\[
x=(x_1,\ldots,x_n),\qquad y=(y_1,\ldots,y_n)\in K^n
\]
y $\alpha,\beta\in K$, entonces
\[
\alpha x+\beta y
=
(\alpha x_1+\beta y_1,\ldots,\alpha x_n+\beta y_n).
\]
Esta expresión será una de las operaciones fundamentales del álgebra lineal y posteriormente recibirá el nombre de combinación lineal de $x$ e $y$.
\end{example}
\begin{proposition}
Sea $K$ un campo. Para cualesquiera $x,y,z\in K^n$ y $\alpha,\beta\in K$ se cumplen
\begin{enumerate}
\item[a)] $x+y=y+x$.
\item[b)] $(x+y)+z=x+(y+z)$.
\item[c)] Existe $0=(0,\ldots,0)\in K^n$ tal que $x+0=x$.
\item[d)] Para cada $x=(x_1,\ldots,x_n)$ existe
\[
-x=(-x_1,\ldots,-x_n)
\]
tal que $x+(-x)=0$.
\item[e)] $\alpha(x+y)=\alpha x+\alpha y$.
\item[f)] $(\alpha+\beta)x=\alpha x+\beta x$.
\item[g)] $(\alpha\beta)x=\alpha(\beta x)$.
\item[h)] $1x=x$.
\end{enumerate}
\end{proposition}
\begin{proof}
Todas las identidades se verifican componente a componente utilizando únicamente los axiomas del campo $K$. Por ejemplo,
\[
\alpha(x+y)
=
\alpha(x_1+y_1,\ldots,x_n+y_n)
\]
\[
=
(\alpha(x_1+y_1),\ldots,\alpha(x_n+y_n))
\]
\[
=
(\alpha x_1+\alpha y_1,\ldots,\alpha x_n+\alpha y_n)
=
\alpha x+\alpha y.
\]
Las demás propiedades se demuestran de manera análoga.
\end{proof}

\section{Definición de espacio vectorial}
Axiomas de suma de vectores y multiplicación por escalares sobre un campo.

\section{Ejemplos fundamentales de espacios vectoriales}
Espacios coordenados, polinomios, matrices, funciones y sucesiones.

\section{Subespacios vectoriales}
Criterios de subespacio, ejemplos y contraejemplos.

\section{Intersección y suma de subespacios}
Operaciones entre subespacios, suma directa y primeros ejemplos de descomposición.
